\section{Methodology}

Building on our observation that compiler errors fall outside typical LLM training distributions, we design a systematic approach to study error format as an independent variable. Our methodology has three components: (1) structured error formatting with controlled conditions, (2) an information preservation test to ensure fair comparison, and (3) a fix specificity framework to investigate optimal hint levels.

\subsection{Overview}

The core idea is to transform raw compiler output into structured templates that align with natural language patterns. Rather than terse, technical tracebacks, we organize error information into labeled sections: PROBLEM, LOCATION, CONTEXT, and ROOT\_CAUSE. Critically, we design multiple format conditions that control for potential confounds:

\begin{table}[H]
\centering
\begin{tabular}{@{}lll@{}}
\toprule
\textbf{Condition} & \textbf{Description} & \textbf{Purpose} \\
\midrule
Raw & Direct compiler/linter output & Baseline \\
Verbose-Raw & Expanded raw output, unstructured & Control for length \\
Structured & Template format with section labels & Treatment \\
Scrambled & Structured content, random section order & Control for information \\
\bottomrule
\end{tabular}
\caption{Experimental conditions for format comparison.}
\end{table}

The Scrambled condition is key to causal inference: it contains identical information to Structured but with randomly permuted section order. If Structured outperforms Scrambled, the effect is due to organizational structure, not information content.

\subsection{Structured Error Format}

Our structured format transforms raw errors into a consistent template:

\begin{lstlisting}
=== ERROR REPORT ===

PROBLEM: [error_type]: [error_message]

LOCATION: Line [line_number] in [file_name]

CONTEXT:
[surrounding_code_snippet]
    ^ [pointer_to_error]

ROOT_CAUSE: [explanation_of_why_error_occurred]

FIX_HINT: [optional, level-dependent guidance]
\end{lstlisting}

\textbf{Design Rationale}: The section labels create parsing anchors that help the model identify relevant information. The consistent structure enables pattern matching against similar examples in training data. The explicit ROOT\_CAUSE section provides causal explanation rather than just symptom description.

\subsubsection{Error Parser Implementation}

We implement a parser that extracts structured information from Python tracebacks:

\begin{lstlisting}[language=Python]
@dataclass
class StructuredError:
    line_number: int
    error_type: str
    error_message: str
    code_context: str
    root_cause: Optional[str] = None
\end{lstlisting}

The parser uses regex-based extraction to identify line numbers, exception types, and relevant code context from raw traceback output. We support 8 common error types: \texttt{NameError}, \texttt{TypeError}, \texttt{IndexError}, \texttt{KeyError}, \texttt{ZeroDivisionError}, \texttt{ValueError}, \texttt{AttributeError}, and \texttt{SyntaxError}.

\subsection{Information Preservation Test}

To ensure our format transformation preserves all diagnostic information, we design a reconstruction test (sub-hypothesis h-m1):

\begin{enumerate}
    \item \textbf{Generate pairs}: Create (raw\_error, structured\_error) pairs for 500 samples across 8 error types
    \item \textbf{Reconstruction task}: Ask a third-party LLM to reconstruct original error details from the structured format
    \item \textbf{Measure accuracy}: Compare reconstructed fields against ground truth
\end{enumerate}

\textbf{Success criterion}: $>$95\% reconstruction accuracy across all fields (line\_number, error\_type, error\_message, code\_context).

This test validates that any performance differences between conditions are due to format, not information loss. Our experiments achieved 100\% reconstruction accuracy in validation mode, confirming information preservation.

\subsection{Fix Specificity Framework}

Motivated by scaffolding theory from HCI research \citep{vygotsky1978zpd}, we investigate whether intermediate-level hints outperform both extremes. We define four specificity levels:

\begin{table}[H]
\centering
\begin{tabular}{@{}lll@{}}
\toprule
\textbf{Level} & \textbf{Name} & \textbf{Example} \\
\midrule
0 & No hint & (diagnostic only) \\
1 & General strategy & ``convert types to match'' \\
2 & Specific pattern & ``use str(x) or int(y)'' \\
3 & Exact fix & ``str(count)'' \\
\bottomrule
\end{tabular}
\caption{Fix specificity levels.}
\end{table}

\textbf{Theoretical basis}: Vygotsky's Zone of Proximal Development suggests that optimal scaffolding provides enough guidance to activate relevant knowledge without bypassing the learning/reasoning process. Level 3 (exact fix) may create copy-paste dependency, while Level 0 provides insufficient guidance.

\subsubsection{Hint Generation}

We implement a hint generator that produces appropriately-leveled fix suggestions:

\begin{lstlisting}[language=Python]
def generate_hint(error: StructuredError, level: int) -> str:
    if level == 0:
        return ""  # No hint
    elif level == 1:
        return get_general_strategy(error.error_type)
    elif level == 2:
        return get_specific_pattern(error)
    else:  # level == 3
        return get_exact_fix(error)
\end{lstlisting}

Hints are generated based on error type and context, with Level 2 providing common fix patterns and Level 3 providing exact code replacements.

\subsection{Experimental Design}

\subsubsection{Format Comparison (h-m2)}

We use a paired comparison design to isolate structure effects:

\begin{enumerate}
    \item \textbf{Sample selection}: 500 code errors from HumanEval+/MBPP+ model failures
    \item \textbf{Condition assignment}: Each error is formatted in both Structured and Scrambled conditions
    \item \textbf{Repair execution}: Apply self-repair with CodeLlama-7B for each condition
    \item \textbf{Outcome measurement}: Binary success/failure per sample
    \item \textbf{Statistical test}: McNemar's test for paired binary outcomes
\end{enumerate}

\textbf{Causal inference}: If Structured $>$ Scrambled with $p < 0.05$, we conclude that organizational structure drives improvement, controlling for information content.

\subsubsection{Fix Specificity Test (h-m3)}

We extend the design to test fix specificity levels:

\begin{enumerate}
    \item \textbf{Within-subject design}: Same error instances across all 4 levels
    \item \textbf{Seeded randomization}: Ensure reproducibility
    \item \textbf{Statistical analysis}: Mixed-effects model with quadratic contrast
\end{enumerate}

\textbf{Prediction}: Inverted-U pattern with peak at Level 1-2.

\subsubsection{Scale Interaction Test (h-c1)}

We test whether format benefits vary by model scale:

\begin{enumerate}
    \item \textbf{Factorial design}: 2 (Format) $\times$ 3 (Model: 7B, 34B, GPT-4)
    \item \textbf{Analysis}: Two-way ANOVA with interaction term
    \item \textbf{Success criterion}: Significant interaction ($p < 0.05$, $\eta^2 > 0.01$)
\end{enumerate}

\subsection{Implementation}

All experiments use the EvalPlus benchmark framework \citep{guo2024evalplus} for evaluation, ensuring rigorous testing with 80$\times$ more test cases than original HumanEval/MBPP. Models are accessed via HuggingFace (CodeLlama) and OpenAI API (GPT-4). Temperature is fixed at 0.0 for deterministic comparison. Maximum repair iterations is set to 5.

Code is organized into modules: \texttt{errors.py} (parsing), \texttt{prompts.py} (formatting), \texttt{models.py} (inference), \texttt{repair\_loop.py} (orchestration), \texttt{evaluate.py} (metrics), and \texttt{analysis.py} (statistics). All code is available in our supplementary materials.
